% intro.tex

%%%%%%%%%%%%%%%%%%%%
\begin{frame}{}
  \begin{exampleblock}{如何使用命题逻辑表达下列命题与推理}
    \begin{center}
      \blue{亚里士多德的``三段论''}

      \[
        \ndnoname{\text{每个人都是要死的 \qquad 苏格拉底是人}}{\text{苏格拉底是要死的}}
      \]
    \end{center}
  \end{exampleblock}

  \pause
  \vspace{0.30cm}
  \[
    \ndnoname{P \qquad Q}{R}
  \]

  \pause
  \vspace{0.50cm}
  \begin{center}
    \red{\bf 命题逻辑无法表达: 个体、全体以及它们之间的关系} \\[15pt]
    我们需要使用表达能力更强的\blue{\bf 一阶谓词逻辑}
  \end{center}
\end{frame}
%%%%%%%%%%%%%%%%%%%%

%%%%%%%%%%%%%%%%%%%%
\begin{frame}{}
  \begin{exampleblock}{例子: 谓词 (Predicate)}
    \begin{enumerate}[<+->][(1)]
      \setlength{\itemsep}{6pt}
      \item 张三是个法外狂徒
      \item 李四与张三是好朋友
      \item 李四也是一个法外狂徒
      \item 王五站在张三与李四中间 (害怕极了)
      \item 王五长得比张三与李四都高
    \end{enumerate}
  \end{exampleblock}

  \pause
  \vspace{0.30cm}
  \begin{center}
    \blue{一元谓词}表达了个体的性质 \\[8pt]
    \pause
    \blue{多元谓词}表达了个体之间的关系 \\[8pt]
    \pause
    \blue{零元谓词}即命题逻辑中的命题符号
  \end{center}
\end{frame}
%%%%%%%%%%%%%%%%%%%%

%%%%%%%%%%%%%%%%%%%%
\begin{frame}{}
  \begin{center}
    一阶谓词逻辑是一个\red{\bf 公理系统}, \\[3pt]
    包含 \teal{\bf 语法 (Syntax)} 与 \teal{\bf 语义 (Semantics)} 两部分。

    \vspace{0.30cm}
    \fig{width = 0.60\textwidth}{figs/syntax-semantics}
    \vspace{0.30cm}

    \pause
    语法与语义是\red{\bf ``对立统一''}的

    \pause
    \vspace{0.50cm}
    {\large \blue{\bf 螺旋上升:} 先讲语法, 再讲语义, 再回到语法, 最后将二者统一起来}
  \end{center}
\end{frame}
%%%%%%%%%%%%%%%%%%%%

%%%%%%%%%%%%%%%%%%%%
% \begin{frame}{}
%   \begin{center}
%     \blue{\large \bf 一阶谓词逻辑}
%     \fig{width = 0.40\textwidth}{figs/Frege}
%     \teal{Gottlob Frege (1848 $\sim$ 1925)}
%   \end{center}
% \end{frame}
%%%%%%%%%%%%%%%%%%%%